% Créditos de las fotografías y de las ilustraciones generadas por IA del
% libro 2 (Química universitaria, primer año), edición española. Los
% registros de licencia legibles por máquina están en
% images/book2/CREDITS.md; mantener ambos archivos sincronizados con la
% versión inglesa (frontmatter/image-credits-book2.tex).
\chapter*{Créditos de las imágenes}

Los dibujos de este libro son originales. Las fotografías, cuando se
usan, se emplean bajo sus respectivas licencias libres, con nuestro
agradecimiento a sus autores; los enlaces completos a las fuentes y las
URL de las licencias de cada fotografía figuran en
\texttt{images/book2/CREDITS.md}, en el repositorio del libro.
Los pictogramas de peligro son los del Sistema Globalmente Armonizado de
clasificación y etiquetado de productos químicos, tal como los dibuja el
paquete \texttt{ghsystem} de \TeX\ Live.

\bigskip
Junto a las fotografías y a los dibujos originales, algunas figuras son
ilustraciones generadas por IA, producidas con la generación de imágenes
de OpenAI a partir de instrucciones escritas por los editores del libro,
y revisadas en cuanto a su exactitud química antes de incluirlas. La
instrucción completa de cada imagen generada está registrada en
\texttt{images/book2/ai/PROMPTS.md}, en el repositorio.

\bigskip
\noindent Fotografías, por capítulo:
\begin{itemize}\small
  \item Cap.~1: Niels Bohr, hacia 1922, AB Lagrelius \& Westphal, dominio
        público (Wikimedia Commons).
  \item Cap.~2: Linus Pauling, 1962, Fundación Nobel, dominio público
        (Wikimedia Commons).
  \item Cap.~5: ejemplar de cobre nativo, de Flipin, OG, CC0 (Wikimedia
        Commons).
  \item Cap.~6: cristales de halita, fluorita y un diamante octaédrico, de
        James St. John, CC BY 2.0; grafito, de Zbynek Burival,
        CC BY-SA 4.0; cristales de nieve de Wilson A. Bentley (1902),
        dominio público (todas de Wikimedia Commons).
  \item Cap.~8: Svante Arrhenius, 1909, Meisenbach Riffarth \& Co.,
        dominio público (Wikimedia Commons).
  \item Cap.~11: precipitados de cloruro, bromuro y yoduro de plata, de
        Rrausch1974, CC BY-SA 3.0 (Wikimedia Commons).
  \item Cap.~12: suspensión de hidróxido de cobre(II), de Alvy16,
        CC BY 4.0; disolución amoniacal de cobre, de Chemicalinterest,
        dominio público (ambas de Wikimedia Commons).
  \item Cap.~13: Walther Nernst, 1889, Smithsonian Libraries, dominio
        público (Wikimedia Commons).
  \item Cap.~16: Jacobus Henricus van 't Hoff, fotografía publicada en
        1911, dominio público (Wikimedia Commons).
  \item Cap.~21: Victor Grignard, dominio público (Wikimedia Commons).
  \item Cap.~25: Vladímir Markóvnikov, retrato procedente de una nota
        necrológica publicada en 1905, dominio público (Wikimedia
        Commons).
  \item Cap.~26: piscinas de litio del Salar de Atacama, NASA Earth
        Observatory, dominio público; sodio en aceite mineral, de Justin
        Urgitis, CC BY-SA 2.5; colores de llama, de Hegelrast,
        CC BY-SA 4.0 (todas de Wikimedia Commons).
  \item Cap.~27: bauxita pisolítica, de James St. John, CC BY 2.0;
        bórax, de Aram Dulyan, dominio público; cuarzo, de Stephanie
        Clifford, CC BY 2.0; fósforo blanco, Hi-Res Images of Chemical
        Elements, CC BY 3.0; Fritz Haber, Fundación Nobel, dominio público
        (todas de Wikimedia Commons).
  \item Cap.~28: extracción de azufre en el cráter del Kawah Ijen, de
        S\'emhur, CC BY-SA 4.0; cloro en una ampolla, de W. Oelen,
        CC BY-SA 3.0; bromo en una ampolla, de Jurii, CC BY 3.0; cristales
        de yodo, de Dnn87, CC BY 3.0 (todas de Wikimedia Commons).
  \item Cap.~29: banco calefactor (banco Kofler), de HBR, CC BY-SA 3.0;
        refractómetro de Abbe, de Foreade, CC BY-SA 4.0 (ambas de
        Wikimedia Commons).
\end{itemize}
\clearpage
